\documentclass[10pt,a4paper]{amsart}
\usepackage[margin=19mm]{geometry}
\usepackage{amsmath,amssymb,mathtools}
\usepackage[scr=rsfs,cal=euler]{mathalfa}
\usepackage{xfrac}
\usepackage[unicode,hidelinks,bookmarks=false]{hyperref}
\newcommand{\ie}{i.e.\ }
\newcommand{\cf}{cf.\ }
\newcommand{\resp}{respectively\ }
\newcommand{\Subsection}{\subsection}
\providecommand{\nobreakdash}{\nobreak\hbox{-}}
\setlength{\emergencystretch}{2em}
\multlinegap0pt
\usepackage{bm}
\usepackage{bbm}
\usepackage{dsfont}
\usepackage{xspace}
\usepackage{cancel}
\usepackage{mathtools}
\usepackage{cleveref}
\usepackage{amsthm}
\makeatletter
\@ifundefined{lemma}{\newtheorem{lemma}{Lemma}}{}
\makeatother
\newcommand{\Id}{\operatorname{Id}}
\newcommand{\oH}{{\bf H}}
\newcommand{\oL}{{\bf L}}
\newcommand{\oP}{{\bf P}}
\newcommand{\oa}{{\bf a}}
\newcommand{\orho}{{\bf \rho}}
\newcommand{\xH}{\mathcal{H}}
\newcommand{\xL}{\mathcal{L}}
\newcommand{\xW}{\mathcal{W}}
\title[Convergence Analysis of Galerkin Approximations for the Lind]{Convergence Analysis of Galerkin Approximations for the Lindblad Master Equation}
\author{R\'emi Robin, Pierre Rouchon}
\date{Source: arXiv:2510.11416v2 (CC BY 4.0)}
\begin{document}
\maketitle

\noindent\emph{Source and licence.} Convergence Analysis of Galerkin Approximations for the Lindblad Master Equation. Version arXiv:2510.11416v2 (CC BY 4.0), distributed under the Creative Commons Attribution 4.0 International licence, as recorded in the arXiv licence metadata for this exact version and shown on the abstract page https://arxiv.org/abs/2510.11416v2. Full rights notice: licenses/arxiv-2510.11416-CC-BY-4.0.txt.\par\smallskip

\noindent\textit{Editorial scope.} Counted unit: the Lemma whose source label is \texttt{lem:key\_estimate} in the article (source0.tex lines 318--322, source label \texttt{lem:key\_estimate}), delivered together with the article's own complete original proof of it (source0.tex lines 324--374, one \texttt{proof} environment of 51 lines). No mathematics is written, repaired or re-derived: statement and proof are the article's own text line for line. Required context retained uncounted: lem:reg\_vs\_rate (lines 290--296). Every definition, display and label of those lines is retained unchanged. Omitted from this package and not redistributed: 1--289, 297--317, 323--323, 375--482. Nothing inside a retained excerpt is omitted or altered: every retained body is delivered exactly as the article's lines are written, so no omission receipt of any kind accompanies this package. Citations: the retained excerpts carry no citation command at all, so no bibliography entry of the article is transcribed and no reference is redistributed. Reference adaptation: none was needed and none was made; every reference inside the delivered unit resolves inside it, so no unresolved reference or citation is produced. Printed slips kept literal: no printed word, symbol, hypothesis or number of the counted statement or of its proof was altered. Layout and class: the article's own class is \texttt{article} with packages including geometry, amssymb, amsmath, amsthm, color, graphicx, dsfont, mathrsfs, stmaryrd, mathtools, braket, authblk, hyperref, cleveref; the wrapper is the lane's pinned \texttt{amsart} editorial wrapper at 10pt on A4, which restates the theorem-like environments the retained text needs and adds the article's own macro definitions that the retained text needs (\texttt{\textbackslash Id}, \texttt{\textbackslash oH}, \texttt{\textbackslash oL}, \texttt{\textbackslash oP}, \texttt{\textbackslash oa}, \texttt{\textbackslash orho}, \texttt{\textbackslash xH}, \texttt{\textbackslash xL}, \texttt{\textbackslash xW}), with extra wrapper packages (bm, bbm, dsfont, xspace, cancel, mathtools, cleveref). The delivered numbering is the unit's own. Assets: no retained span contains a figure environment, a graphics-inclusion command, a PDF-page inclusion, a table or a TikZ picture; no image file is copied into this package, the package declares no asset dependency and no image file is redistributed with it. Licence: version arXiv:2510.11416v2 of this article is distributed under the Creative Commons Attribution 4.0 International licence, as recorded in the arXiv licence metadata for this exact version and shown on the abstract page https://arxiv.org/abs/2510.11416v2. A full rights notice is delivered at licenses/arxiv-2510.11416-CC-BY-4.0.txt. Characters: every retained excerpt is pure ASCII, so the delivered engine, XeLaTeX with its default Latin Modern text font, reports no missing character.
\begin{lemma}
    \label{lem:reg_vs_rate}
    Let $0 \leq s_1 \leq s_2$ and consider $\oP_N$ the spectral projector on $[0,N]$ of $\Id + \oa^\dag \oa$. The following inequality holds
    \begin{align}
        \|\oP_N^\perp\|_{\xH^{s_2}\to \xH^{s_1}}\leq \frac{1}{N^{(s_2-s_1)/2}}.
    \end{align}
\end{lemma}

\begin{lemma}
    \label{lem:key_estimate}
    There exists $C_k>0$ such that for all $\orho\in \xW^{k,1}$,
    $\| (\xL-\xL_N)(\orho)\|_1\leq \frac{C_k}{N^{(k-d)/2}} \|\orho\|_{\xW^{k,1}}$
\end{lemma}

\begin{proof}
    We establish the inequality for $\orho\in \xW^f$, a simple density argument then allows us to extend to $\xW^{k,1}$. Let us start with the Hamiltonian part of the Lindbladian:
    \begin{align}
        \|[\oH-\oH_N,\orho]\|_1 &\leq 2 \| (\oH-\oH_N)\orho\|_1\\
        & \leq 2 (\| \oP_N \oH \oP_N^\perp \orho \|_1 + \| \oP_N^\perp \oH \orho \|_1)\\
        & \leq 2 (\| \oH \oP_N^\perp \orho \|_1 + \| \oP_N^\perp \oH \orho \|_1).
    \end{align}
    Then, using \cref{lem:reg_vs_rate}, we get
    \begin{align}
        \| \oH \oP_N^\perp \orho \|_1+ \| \oP_N^\perp \oH \orho \|_1&
        \leq \left( \|\oH\|_{\xH^{p_H}\to \xH} \|\oP_N^\perp\|_{\xH^k \to \xH^{p_H}} + \|\oP_N^\perp\|_{\xH^{k-p_H}\to \xH} \|\oH\|_{\xH^{k}\to \xH^{k-p_H}} \right)\|\orho\|_{\xW^{k,1}}\\
        &\leq \frac{\|\oH\|_{\xH^{p_H}\to \xH}+\|\oH\|_{\xH^{k}\to \xH^{k-p_H}}}{N^{(k-p_H)/2}} \|\orho\|_{\xW^{k,1}}.
    \end{align}

    Let us now address the dissipative part, for a jump operator $\oL$, we have
    \begin{align}
        \label{eq_dissipator_main}
        \|(D[\oL]-D[\oL_N])(\orho)\|_1 \leq \|\oL \orho \oL^\dag - \oL_N\orho \oL_N^\dag\|_1 + 2\|(\oL^\dag \oL-\oL_N^\dag \oL_N)\orho\|_1
    \end{align}
    \paragraph{First term of the right-hand side of \cref{eq_dissipator_main}.} We split again:
    \begin{align}
        \label{eq_second_decompo}
        \|\oL \orho \oL^\dag -\oL_N\orho \oL_N^\dag\|_1\leq \|(\oL-\oL_N)\orho \oL^\dag\|_1 +\|\oL_N\orho (\oL^\dag-\oL_N^\dag)\|_1
    \end{align}
    The first term is handled using $ \oL -\oL_N= \oP_N^\perp \oL + \oP_N \oL \oP_N^\perp$, namely:
    \begin{align}
        \|(\oL-\oL_N)\orho \oL^\dag\|_1 & \leq \|\oP_N^\perp \oL \orho \oL^\dag\|_1+ \|\oP_N \oL \oP_N^\perp \orho \oL^\dag\|_1\\
        & \leq \frac{1}{N^{(k-p_j)/2}}\left( \|\oL \|_{\xH^{p_j}\to \xH}+\|\oL \|_{\xH^{k}\to \xH^{k-p_j}} \right) \|\orho\|_{\xW^{k,1}}\|\oL \|_{\xH^{p_j}\to \xH}.
    \end{align}
    Similarly, for the second term of \cref{eq_second_decompo}, we have
    \begin{align}
        \|\oL_N\orho (\oL^\dag- \oL_N^\dag)\|_1 &\leq \frac{1}{N^{(k-p_j)/2}}\|\oL_N\|_{\xH^{p_j}\to \xH} \|\orho\|_{\xW^{k,1}}\left(\|\oL \|_{\xH^{p_j}\to \xH}+\|\oL \|_{\xH^{k-p_j}\to \xH^{p_j}}\right) \\
        &\leq \frac{1}{N^{(k-p_j)/2}}\|\oL \|_{\xH^{p_j}\to \xH} \|\orho\|_{\xW^{k,1}} \left(\|\oL \|_{\xH^{p_j}\to \xH}+\|\oL \|_{\xH^{k-p_j}\to \xH^{p_j}}\right).
    \end{align}
    \paragraph{Second term of the right-hand side of \cref{eq_dissipator_main}.}
    We decompose the left factor as
    \begin{align}
        \label{eq_remain_second_part}
        (\oL^\dag \oL - \oP_N \oL^\dag \oP_N \oL \oP_N)\orho= (\oP_N^\perp \oL^\dag \oL + \oP_N \oL^\dag \oP_N^\perp \oL+ \oP_N \oL^\dag \oP_N \oL \oP_N^\perp)\orho
    \end{align}
    For the first term of the sum, we get
    \begin{align}
        \|\oP_N^\perp \oL^\dag \oL \orho \|_1&\leq \|\oP_N^\perp\|_{\xH^{k-2p_j}\to \xH} \| \oL^\dag \oL\|_{\xH^{k}\to \xH^{k-2p_j}} \|\orho\|_{\xW^{k,1}}\\
        &\leq \frac{1}{N^{(k-2p_j)/2}}\| \oL^\dag \oL\|_{\xH^{k}\to \xH^{k-2p_j}} \|\orho\|_{\xW^{k,1}},
    \end{align}
    The remaining term of \cref{eq_remain_second_part} are treated the same way.
    Hence, we have shown that
    \begin{align}
        \| (\xL-\xL_N)(\orho)\|_1\leq \frac{C_k}{N^{(k-\max (p_H,2p_j))/2}} \|\orho\|_{\xW^{k,1}}
    \end{align}
\end{proof}

\end{document}
